% GENERATED by tools/build_m05_code_sheet.py -- do not edit.
\begin{codebox}{Box 3}{The stochastic block model, with graph-tool}
import graph_tool.all as gt
import numpy as np
gt.seed_rng(1)
g = gt.Graph(directed=False)
g.add_edge_list(karate.get_edgelist())
state = gt.minimize_blockmodel_dl(g, state_args=dict(deg_corr=False))
sbm = state.get_blocks().a
print(sbm)
groups = np.unique(sbm, return_inverse=True)[1]
colors = g.new_vertex_property("vector<double>")
for v in g.vertices():
    colors[v] = (*palette[groups[int(v)]], 1)
pos = gt.graph_draw(
    g,
    vertex_fill_color=colors,
    vertex_text=g.vertex_index,
)
\end{codebox}
\begin{linenotes}
\item[1] graph-tool is a second library. It fits the stochastic block model, which you know from the lecture. \texttt{as gt} is its short name.
\item[2] \texttt{numpy} is a library for lists of numbers. \texttt{np} is its short name.
\item[3] The fit makes random choices. \texttt{seed\_rng(1)} fixes them, so that everyone in the room gets the same answer. Change the \texttt{1} and you may get a different grouping.
\item[4] graph-tool cannot read igraph's networks. It has its own kind, so we start with an empty one. \texttt{directed=False} says that a friendship goes both ways.
\item[5] \texttt{karate.get\_edgelist()} asks igraph for all the friendships as pairs of node numbers. \texttt{add\_edge\_\allowbreak list} hands the pairs to graph-tool, which creates the nodes and the edges. This line is the only bridge between the two libraries.
\item[6] \texttt{minimize\_blockmodel\_dl} fits the stochastic block model. \texttt{dl} stands for \emph{description length}: of all the ways to group the nodes, it picks the one that describes the network in the fewest bits. It decides the number of groups too. \texttt{state\_args=dict(deg\_corr=False)} asks for the plain model of the lecture, where the group alone decides how likely a friendship is. Graph-tool's default also lets every member have a number of friends of their own; on a network this small, that version finds one single group. The fitted model is stored as \texttt{state}.
\item[7] \texttt{state.get\_blocks()} gives the group of each node, in graph-tool's own array type. \texttt{.a} turns it into a plain array of numbers, like the list in Box 2.
\item[8] Print it: again one number for each of the 34 nodes.
\item[9] graph-tool numbers its groups with gaps (say 15 and 32). \texttt{np.unique} with \texttt{return\_\allowbreak inverse=True} renumbers them 0, 1, 2, \ldots{} so that they fit the palette; the \texttt{[1]} picks that renumbered list.
\item[10] graph-tool keeps what it knows about nodes in a \emph{property map}: a table with one value per node. This one will hold a list (red, green, blue) for each node.
\item[11--12] A loop over the nodes \texttt{v} of the network. \texttt{groups[int(v)]} is the group of node \texttt{v} (\texttt{int} turns graph-tool's node object into its number), \texttt{palette[\ldots]} is the colour of that group, and \texttt{(*palette[\ldots], 1)} unpacks its red, green and blue and adds a 1 for fully opaque. The result goes into the table.
\item[13--17] \texttt{gt.graph\_draw} is graph-tool's drawing function. \texttt{vertex\_fill\_color=colors} fills each node with its colour from the table, and \texttt{vertex\_text=g.vertex\_index} writes its number on it. The function hands back the position it gave each node; we keep that as \texttt{pos}, which also stops Colab from printing it under the picture.
\end{linenotes}
